\section{第三阶段深讲：特权边界、装载和内存提交点}

系统调用、exec 和缺页处理都在改变“谁能访问什么”。这类代码最怕半提交：权限已经放开但对象还没准备好，地址空间已经切换但旧程序还要返回错误，页表已经失效但物理页提前释放。学习这一章时，应把每条路径拆成准备、验证、提交、回滚四个阶段，而不是只记住 API 名字。

\subsection{边界代码的共同结构}

边界代码有三个共同问题。第一，输入来自不可信主体：系统调用参数来自用户寄存器和用户内存，ELF 文件可能被攻击者构造，page fault 地址可能来自错误或恶意访问。第二，执行路径会改变全局可见状态：fd 表、地址空间、页表、文件页、凭据和调度状态。第三，失败必须可解释：有些失败返回错误码，有些发送信号，有些说明内核自身已经损坏。

\begin{lstlisting}
boundary operation:
  snapshot untrusted input
  validate object identity and permissions
  allocate temporary resources
  build new state off to the side
  commit with a short ordered sequence
  publish visibility only after invariants hold
  roll back temporary resources on every pre-commit failure
\end{lstlisting}

这个结构适用于 \code{openat}、\code{execve}、\code{mmap}、\code{mprotect}、COW fault 和文件写回。源码阅读时，只要找到“提交点”在哪里，前后语义就会清楚很多：提交点之前失败要回到旧状态，提交点之后失败只能按新状态处理或终止任务。

\subsection{系统调用返回值也是 ABI}

系统调用的结果不只是一个整数。ABI 规定返回值寄存器、错误码编码、哪些寄存器被破坏、哪些状态在返回前必须处理。libc wrapper 再把内核返回值转换成 \code{-1} 和 \code{errno}。如果教学 OS 忽略这些细节，用户程序就只能靠猜。

\begin{longtable}{p{0.22\linewidth}p{0.34\linewidth}p{0.32\linewidth}}
\toprule
返回类型 & 例子 & 调用者应该怎样理解 \\
\midrule
非负数量 & \code{read/write} 返回字节数 & 可能短读/短写，不等于请求全部完成 \\
新句柄 & \code{open/socket/dup} 返回 fd & fd 是进程表索引，不是对象永久身份 \\
0 表示成功 & \code{close/mprotect/munmap} & 只说明接口语义完成，不说明持久化 \\
\code{-errno} & 内核内部常见约定 & wrapper 转成 \code{errno} 或直接传递 \\
信号 & 用户访存错误、非法指令 & 错误发生在用户执行流，不一定有 syscall 返回 \\
\bottomrule
\end{longtable}

\code{close} 的错误尤其容易被误解。fd 表项可能已经释放并可被复用，即使底层写回错误在 close 期间才报告。用户态不能在 \code{close} 返回错误后再重试同一个 fd，因为这个整数可能已经指向别的文件。正确设计要把持久化错误尽量通过 \code{fsync} 这类明确边界暴露。

\subsection{restartable syscall 和 EINTR}

阻塞系统调用可能在等待期间收到信号。内核有时返回 \code{EINTR}，有时把系统调用改写成可重启状态，让信号处理返回后重新执行。重启不是简单“再调一次函数”，因为部分 I/O、超时参数和用户可见副作用必须精确处理。

\begin{lstlisting}
blocking read:
  copy no bytes yet
  sleep waiting for data
  signal arrives
  if handler installed:
    either return -EINTR
    or arrange restart after signal frame

partial read:
  copied some bytes
  signal arrives
  return copied bytes, not -EINTR
\end{lstlisting}

这个规则解释了为什么系统调用实现要记录“是否已经产生用户可见进展”。已经复制给用户的字节、已经消耗的 socket 数据、已经创建的 fd，不能因为信号到来就假装从未发生。测试时应同时覆盖“阻塞前被打断”和“部分完成后被打断”。

\subsection{copy\_from\_user 的分段提交}

大 \code{read/write} 不应一次性申请巨大内核缓冲。常见实现会分段处理用户 iovec 或页片段。分段处理带来短完成语义：前几段成功，后面遇到坏页或信号，返回值可能是已经完成的字节数。

\begin{lstlisting}
write_iter(file, iov):
  done = 0
  for each segment in iov:
    while segment has bytes:
      chunk = copy_chunk_from_user(segment)
      if chunk fault:
        return done > 0 ? done : -EFAULT
      ret = file_write_chunk(file, chunk)
      if ret < 0:
        return done > 0 ? done : ret
      done += ret
  return done
\end{lstlisting}

这里的提交粒度是 chunk。若 chunk 已写入 pipe 或 socket queue，就不能在后续 fault 时撤销。文件系统写入还要区分 page cache dirty、journal commit 和设备持久化。一本教材若只说“复制用户 buffer 再写文件”，读者会漏掉短写、EFAULT 和持久化边界。

\subsection{exec 的三道权限门}

\code{execve} 至少经过三道权限门。第一道是路径和文件权限：调用者是否能走到该文件，文件是否可执行，挂载是否 \code{noexec}。第二道是二进制格式和装载权限：ELF 段权限是否合理，解释器是否允许执行，栈是否可执行。第三道是凭据转换：setuid、file capability、LSM、no\_new\_privs 和 ptrace 规则共同决定新程序的身份。

\begin{lstlisting}
exec permission gates:
  path permission and executable inode
  binary format and memory layout policy
  credential transition and security hooks
\end{lstlisting}

这三道门的错误码和失败时机不同。路径不存在是 \code{ENOENT}；不是可执行格式可能是 \code{ENOEXEC}；参数过大是 \code{E2BIG}；权限不足是 \code{EACCES} 或 \code{EPERM}；解释器不存在也会让整个 exec 失败，但旧程序必须还能拿到错误。把这些都归为“启动失败”，无法写出 shell、容器运行时或服务管理器。

\subsection{ELF 段权限不是 section 权限}

运行时装载看 program header，链接和调试看 section header。很多 ELF 被 strip 后没有完整 section 信息，仍然能运行。教学装载器如果读取 \code{.text/.data/.bss} 名字来建立映射，就是把链接器视角误当成内核视角。

\begin{longtable}{p{0.22\linewidth}p{0.34\linewidth}p{0.32\linewidth}}
\toprule
视角 & 使用对象 & 典型消费者 \\
\midrule
运行时装载 & \code{PT\_LOAD/PT\_INTERP/PT\_PHDR} & 内核 exec、动态链接器 \\
链接和重定位 & section、symbol、relocation & 链接器、objdump、调试器 \\
调试展开 & DWARF、CFI、eh\_frame & gdb、异常展开、profiler \\
\bottomrule
\end{longtable}

段映射还要处理文件尾页。若 \code{p\_filesz} 结束在页中间，而 \code{p\_memsz} 更大，同一页中从文件末尾到页末的部分必须清零；后续整页 BSS 也要清零或懒分配零页。这个细节若错，未初始化全局变量会读到旧文件内容或未清理内存。

\subsection{地址空间切换的提交顺序}

exec 提交时，最危险的是先激活新页表再访问旧地址空间数据，或先释放旧 mm 再需要从旧用户栈取参数。安全做法是：提交前把 path、argv、envp、解释器路径和 auxv 所需数据复制到内核或新栈页；新 mm 完整可用后再短序列切换。

\begin{lstlisting}
safe exec commit:
  prepare new_mm, new_stack, new_regs
  copy all old user data before old_mm is dropped
  stop sibling threads if multithreaded
  install credentials that belong to new image
  install current.mm = new_mm
  switch hardware page table
  make old_mm unreachable from current
  drop old_mm references
  return to new user entry
\end{lstlisting}

顺序里有两个可见性点：\code{current.mm} 对内核其他路径可见，硬件页表根对 CPU 访存可见。多核系统还要考虑同一地址空间是否仍在其他 CPU 活跃。若旧 mm 的页表页在远端 CPU TLB 失效前被释放，后果和普通 \code{munmap} 的 shootdown 错误一样严重。

\subsection{mprotect 的权限收缩和扩大}

\code{mprotect} 改权限时，扩大权限和收缩权限的风险不同。扩大权限要检查 VMA 的 \code{may} 权限和文件打开模式；收缩权限要立刻让硬件旧权限失效。比如把 RW 改成 R，如果不 flush TLB，CPU 可能继续使用缓存的 writable 翻译。

\begin{lstlisting}
mprotect(range, new_prot):
  split vmas at range boundaries
  for each vma in range:
    if new_prot exceeds vma.may_prot:
      rollback splits if needed
      return -EACCES
  update vma.prot
  update present PTE permission bits
  flush TLB for affected range
\end{lstlisting}

VMA 切分是另一个回滚点。若先切分出三个 VMA，后续某段权限检查失败，需要合并回原状态或保证中间状态不会被观察到。真实内核会通过锁保护 mm 结构，让其他线程不能在半修改区间中处理 fault。

\subsection{munmap 和延迟释放}

\code{munmap} 删除映射后，用户不能再访问这段虚拟地址。但对应物理页不一定立刻释放：其他进程可能映射同一文件页，I/O 可能持有页引用，RCU 或 TLB gather 可能延迟释放页表页。语义上的不可访问和资源上的最终释放是两件事。

\begin{lstlisting}
munmap(range):
  remove or split VMAs
  clear PTEs in range
  collect old PTE pages and mapped pages into batch
  flush TLB for range on all active CPUs
  drop page references after stale translations are gone
  free page table pages when no walker can see them
\end{lstlisting}

这能解释为什么内存统计不是瞬时下降，也解释为什么页表释放要和 TLB shootdown、页表遍历锁、RCU 等机制配合。教学 OS 可简化，但必须保留顺序原则：先撤销可达性，再失效旧翻译，最后释放底层资源。

\subsection{page fault 的锁顺序}

缺页路径会碰到 mm 锁、VMA 锁、页表锁、page cache 锁、folio 锁、文件系统锁和 I/O 等待。锁顺序错了会出现非常隐蔽的死锁。一个基本原则是：先稳定地址空间和 VMA，再稳定页表项，再稳定物理页或文件页；等待 I/O 时尽量不持有会阻塞其他 fault 的大锁。

\begin{lstlisting}
file page fault:
  take mmap read lock or VMA lock
  find and validate VMA
  locate page cache folio
  if folio locked or needs I/O:
    drop/reduce locks that block unrelated faults
    wait for folio
    retry validation
  take page table lock
  install PTE if still valid
  drop locks
\end{lstlisting}

这里的 retry 很重要。等待 I/O 期间，另一个线程可能 \code{munmap} 或 \code{mprotect} 了同一区间。醒来后不能直接安装 PTE，必须重新验证 VMA 和权限。许多内存管理 bug 都来自“睡眠前检查过，醒来后不再检查”。

\subsection{源码阅读路线}

这一阶段读 Linux 源码可以按路径而不是文件顺序推进。

\begin{longtable}{p{0.27\linewidth}p{0.38\linewidth}p{0.23\linewidth}}
\toprule
路径 & 关键文件或函数 & 读法 \\
\midrule
系统调用入口 & \code{arch/x86/entry/}、\code{kernel/seccomp.c} & 参数、seccomp、返回工作 \\
用户复制 & \code{arch/x86/lib/usercopy*}、\code{lib/iov\_iter.c} & fault fixup、短拷贝、iov \\
exec/ELF & \code{fs/exec.c}、\code{fs/binfmt\_elf.c} & 准备新镜像和提交点 \\
VMA 管理 & \code{mm/mmap.c}、\code{include/linux/mm\_types.h} & maple tree、切分、合并 \\
缺页 & \code{mm/memory.c}、\code{mm/filemap.c} & anon/file/COW fault 分支 \\
页表失效 & \code{mm/tlb.c}、架构 TLB 代码 & batch、shootdown、延迟释放 \\
回收/OOM & \code{mm/vmscan.c}、\code{mm/oom\_kill.c} & reclaim 顺序和 victim 选择 \\
\bottomrule
\end{longtable}

每条路径都按同一张表做笔记：输入对象、持有锁、可睡眠点、提交点、回滚对象、用户可见结果。这样能把几万行源码压成可理解的状态机。

\subsection{边界反例集}

高质量测试要故意打破边界。

\begin{longtable}{p{0.27\linewidth}p{0.38\linewidth}p{0.23\linewidth}}
\toprule
反例 & 构造方法 & 必须观察到 \\
\midrule
用户指针半页坏 & iovec 前半有效后半 unmapped & 短完成或 \code{EFAULT}，内核不崩 \\
exec 坏解释器 & ELF 的 \code{PT\_INTERP} 指向不存在路径 & exec 失败，旧程序继续 \\
ELF 溢出段 & \code{p\_vaddr + p\_memsz} 整数溢出 & 拒绝装载 \\
\code{mprotect} 后旧权限 & RW 改 R 后立即写 & 必须 fault，不能被 TLB 旧权限放行 \\
fault 等 I/O 时 munmap & 线程 A fault 等页，线程 B unmap & A 醒来后重新验证并失败 \\
close 与 read 并发 & 一个线程 read 已 fget，另一个 close 并复用 fd & read 用旧 file 引用，fd 数字可复用 \\
memcg OOM & cgroup 内触发限制 & 只在 cgroup 内选择 victim \\
\bottomrule
\end{longtable}

这些反例比正常路径更能检验实现质量。正常路径只证明代码能跑，反例才能证明边界没有漏。
